% model.tex — security model for the HMK-1 wrap, and our own cryptanalysis
% of the implementation. Claims here are stated with the assumptions they
% actually use; anything not proved is listed in the limitations section.

\section{Security model for the wrap}
\label{sec:model}

The hybrid-KEM and signature bounds of Sec.~\ref{sec:thm1} say nothing about
the sealing layer itself. This section states what the wrap of
Alg.~\ref{alg:wrap} guarantees on its own, so that the argument does not rest
on the intuition that HMAC is a good primitive. The adversary is a
probabilistic polynomial-time algorithm with access to the public parameters
and to the algorithms of Alg.~\ref{alg:wrap}; the master key $mk$ is a
uniformly random 32-byte string.

\subsection{Games}

\begin{definition}[Sealing games]
\label{def:games}
Let $\Pi$ be the wrap of Alg.~\ref{alg:wrap} under a fixed secret $ss$.
\begin{itemize}
  \item[$\mathsf{IND}$:] $\mathcal{A}$ may query a sealing oracle. On input
    $aad_i$ the challenger draws a fresh nonce $n_i$, draws
    $b \xleftarrow{\$} \{0,1\}$ and returns
    $(ct_i, tag_i) \gets \Pi(ss, n_i, m_b, aad_i)$ for challenge keys
    $m_0, m_1$ chosen by $\mathcal{A}$. $\mathcal{A}$ wins by outputting $b$.
  \item[$\mathsf{CTXT}$:] $\mathcal{A}$ may query the sealing oracle
    adaptively (any AAD, fresh nonces) and wins by outputting
    $(n^\ast, ct^\ast, tag^\ast, aad^\ast)$ that unwraps successfully and was
    never returned by the oracle.
\end{itemize}
\end{definition}

\begin{theorem}[Sealing security]
\label{thm:seal}
Let $F$ be the keyed function HMAC-SHA256 under the derived keys of
Alg.~\ref{alg:wrap}. If $F$ is a secure pseudorandom function, then for every
adversary making $q'$ unwrap attempts,
\[
  \adv{\mathsf{IND}}_{\Pi}(\mathcal{A}) \;\le\; 2\,\adv{\mathrm{PRF}}_{F}(\mathcal{B}),
  \qquad
  \adv{\mathsf{CTXT}}_{\Pi}(\mathcal{A}) \;\le\;
  \adv{\mathrm{PRF}}_{F}(\mathcal{B}') + \frac{q'}{2^{256}} .
\]
\end{theorem}

\begin{proof}[Proof sketch]
\emph{Confidentiality.} Each record's keystream is
$KS_i = F_{kek}(\texttt{"hmk/v1/ks"} \concat n_i)$. With distinct nonces these
are distinct PRF inputs, so a hybrid argument over the $q$ queries replaces
every $KS_i$ by an independent uniform 32-byte string at cost
$\adv{\mathrm{PRF}}_{F}$, after which $ct_i = m_b \oplus KS_i$ is uniform and
independent of $b$; answering with a real keystream instead of a uniform one
costs a second factor of $\adv{\mathrm{PRF}}_{F}$, hence the factor two.

\emph{Integrity.} A tag
$tag^\ast = F_{kmac}(\texttt{"hmk/v1/tag"} \concat n^\ast \concat ct^\ast
\concat aad^\ast)$ for a record the oracle never produced is a PRF guess on an
input not previously evaluated. Up to $\adv{\mathrm{PRF}}_{F}$ it succeeds with
probability at most $2^{-256}$ per attempt, and a union bound over $q'$
attempts gives the stated bound. The reduction needs no collision resistance
from HMAC: the tag input is a fixed-length, prefix-free encoding
(Sec.~\ref{sec:enc}), so distinct records give distinct tag inputs.
\end{proof}

\subsection{Fixed-length encodings remove the parsing surface}
\label{sec:enc}

Every field of a record has exactly one length: $ss=64$, $n=16$, $ct=32$,
$tag=32$ bytes, and the AAD is the only variable-length field, absorbed last
into the tagged message. Consequently (i) every record has exactly one byte
encoding, so two parsers cannot disagree about the same bytes; (ii) there is no
length field to reinterpret, so no truncation or length-extension variant
changes what was authenticated; and (iii) the tagged string is prefix-free,
which is what the forgery argument above needs. The same reasoning is why the
two signature encodings (fixed-width $r \concat s$ for ES256, DER underneath
ML-DSA) are normalised at the boundary in Sec.~\ref{sec:hmk}: a
variable-length encoding inside an authenticated record reintroduces exactly
the ambiguity that (i)--(iii) exclude.

\subsection{Context binding}
\label{sec:aad}

\begin{theorem}[No cross-context replay]
\label{thm:ctx}
If $F$ is a PRF, then for any two distinct contexts $aad \neq aad'$ a record
produced under $aad$ fails to unwrap under $aad'$ except with probability
$2^{-256}$ up to $\adv{\mathrm{PRF}}_{F}$, even if the same nonce and the same
secret are reused.
\end{theorem}

\begin{proof}
The AAD enters only the tag input, so the two tag inputs are
$\texttt{"hmk/v1/tag"} \concat n \concat ct \concat aad$ and the same string
ending in $aad'$. The common prefix has fixed length, so the inputs differ;
PRF outputs on distinct inputs are indistinguishable from independent uniform
values, and coincide with probability $2^{-256}$. The ciphertext is unchanged
because the keystream does not depend on the AAD, which is why binding must
come from the tag --- and why the tag is verified before any key bytes are
released (\textsf{cmov} in Alg.~\ref{alg:wrap}).
\end{proof}

\noindent Theorem~\ref{thm:ctx} is what the deployment uses in place of a
session identifier: a master key sealed for
\texttt{rp=https://\ldots;sub=\ldots;aud=\ldots} is not accepted in another
relying-party context, so a stored record cannot be replayed into a different
tenant. \path{TestNonceReuseIsDetectableAcrossContexts} exercises this case
with a deliberately repeated nonce.

\subsection{Freshness is the one thing the construction cannot enforce}
\label{sec:fresh}

\begin{proposition}[Nonce reuse]
\label{prop:nonce}
Let two records be produced under the same $ss$, the same $aad$ and the same
nonce $n$, but different keys $m, m'$. Then $ct \oplus ct' = m \oplus m'$
exactly: the keystream cancels and either plaintext is recoverable from the
other given one ciphertext.
\end{proposition}

\begin{proof}
$KS = F_{kek}(\texttt{"hmk/v1/ks"} \concat n)$ depends on $(ss, n)$ only, so
$ct \oplus ct' = (m \oplus KS) \oplus (m' \oplus KS) = m \oplus m'$.
\end{proof}

\noindent This is the usual one-time-pad discipline of a nonce-based
keystream. We state it as a proposition rather than leaving it in a doc comment
because it is the design's only caller obligation. The implementation cannot
detect reuse without keeping state; what it does instead is (i) document the
obligation at the API, (ii) derive every nonce in the demo and the TUI from the
OS CSPRNG at wrap time, and (iii) pin the behaviour in
\path{TestNonceReuseIsDetectableAcrossContexts}, so a future attempt to
``optimise'' the nonce into a counter or a constant fails a named test. We
considered deriving the nonce from the master key itself and rejected it: that
makes reuse impossible but turns the wrap into deterministic encryption, which
leaks equality of master keys across tenants.

\subsection{Cryptanalysis of our own implementation}
\label{sec:selfattack}
We attacked the implementation, not just the design. Two findings are worth
recording, both found by reading code and both fixed.

\paragraph{A truncated domain-separation label.}
Alg.~\ref{alg:wrap} uses three labels: \texttt{hmk/v1/rk}, \texttt{hmk/v1/ks}
and \texttt{hmk/v1/tag}. They are 9, 9 and 10 bytes long, but the C core copied
a hard-coded \texttt{9} bytes in both places where \texttt{hmk/v1/tag} is
absorbed, so the tag was in fact computed over \texttt{hmk/v1/ta}. No
cross-protocol collision followed --- the three truncated labels remain
pairwise distinct --- but the separation the specification promises was not the
separation the code implemented, and any future rename (for example to
\texttt{hmk/v1/tag2}) would have collided with the old label silently. The fix
derives every length from \texttt{sizeof} on the label array, so the class of
defect is now a compile-time impossibility rather than something a reviewer has
to notice. This is the sort of bug that tests of \emph{behaviour} never catch:
every round-trip, tamper and constant-time test passed both before and after.

\paragraph{Cross-implementation known-answer vectors.}
Agreement between a C core and its Go bindings is not evidence that the
construction is what the paper says. \path{reference/hmk_ref.py} is an
independent implementation written from the specification in this section
(big-integer field arithmetic, a byte-at-a-time diffusion map, the standard
library's HMAC), sharing no code with \path{csrc/}.\footnote{Its one shared
assumption is SHA-256 itself, which both call from their standard library.}
\path{tools/gen_kat.py} derives six vectors from it --- including degenerate
all-zero inputs, a zero field operand, operands larger than $p$, and an OIDC
context AAD --- and commits them to \path{kat/hmk_vectors.json}.
\path{TestCOREAgreesWithReference} then drives the production cgo path and
requires byte-identical ciphertext \emph{and} tag for every vector, so the two
implementations cannot drift apart without a red test. The same vectors pin the
wire format: a change to a label, a domain prefix or the field prime changes
the vector file, and \texttt{make kat} fails until the change is intentional.

\paragraph{What this does and does not establish.}
The reductions above are standard-model statements about HMAC-SHA256 and say
nothing about the field combiner or the $\F_2$ layer; those are argued
separately (Prop.~\ref{prop:onebreak} and the verified rank/support properties
of $M$) and are not needed for confidentiality, since the keystream already
depends on the full 64-byte secret through the PRF input. They matter as
belt-and-braces: if one half of the hybrid secret is broken, $K_f$ retains the
unbroken half's entropy instead of collapsing to a function of the broken one.
